\section*{Evidence and verification}

The status of a formal result must be reported at the level actually checked.
These levels are not interchangeable:
\begin{enumerate}
  \item a displayed calculation checks only its displayed instances;
  \item a truth table checks all valuations of the finitely many variables it
        contains;
  \item a derivation checks derivability in the stated calculus;
  \item a proof establishes the stated result under its stated
        hypotheses;
  \item a proof-assistant kernel checks a formal proof relative to the kernel,
        imported declarations, and admitted axioms;
  \item successful \LaTeX{} compilation checks typesetting syntax, not logical
        correctness; and
  \item editorial or formal-logical review is independent evidence and must not be
        inferred from any preceding level.
\end{enumerate}

When this document states a result status explicitly, it uses exactly one of
three terms.  \emph{Complete proof} means that all substantive formal-logical
steps are supplied in this document.  \emph{Proof sketch} identifies the
construction and states which substantive verification is omitted.
\emph{Cited theorem} states exact hypotheses and a conclusion but refers to a
named source for the proof.  A dependency used within a complete proof is
stated separately as a cited theorem; it is not hidden in a qualified status.

